\documentclass[10pt,letterpaper]{article}

% =========================================================================
% PAQUETES Y CONFIGURACIÓN TIPOGRÁFICA
% =========================================================================
\usepackage[utf8]{inputenc}
\usepackage[T1]{fontenc}
\usepackage[spanish]{babel} % Por defecto usa "Cuadro" según la norma institucional y el ejemplo
\usepackage{lmodern}        % Tipografía Latin Modern idéntica al documento de referencia
\usepackage{amssymb}        % Símbolos matemáticos y geométricos
\usepackage{indentfirst}    % Sangría en el primer párrafo tras los títulos

% Nota: Si algún evaluador le exige expresamente Times New Roman 12pt según Formato (1),
% descomente la siguiente línea y cambie 10pt por 12pt en \documentclass:
% \usepackage{mathptmx}

% Configuración de márgenes: 2.0 cm en los cuatro bordes según Ejemplo_formato_entregables.pdf
\usepackage[letterpaper,top=2.0cm,bottom=2.0cm,left=2.0cm,right=2.0cm,footskip=1.0cm]{geometry}

\usepackage{graphicx}
\usepackage{array}
\usepackage{tabularx}
\usepackage{multirow}
\usepackage{longtable}
\usepackage{booktabs}       % Tablas de calidad tipográfica formal (toprule, midrule, bottomrule)
\usepackage{colortbl}
\usepackage[table]{xcolor}
\usepackage{fancyhdr}
\usepackage{enumitem}
\usepackage{titlesec}
\usepackage{setspace}
\usepackage{lastpage}
\usepackage[hidelinks]{hyperref}

% =========================================================================
% FORMATO DE PÁRRAFOS Y LISTAS
% =========================================================================
% Viñetas cuadradas rellenas de 3pt x 3pt idénticas al PDF de referencia
\setlist[itemize]{
  label=\rule[0.2ex]{3pt}{3pt},
  leftmargin=1.5em,
  topsep=3pt,
  itemsep=3pt,
  parsep=0pt
}
\setlist[enumerate]{
  leftmargin=1.5em,
  topsep=3pt,
  itemsep=3pt,
  parsep=0pt
}

% Sangría e interlineado idénticos al documento de referencia
\setlength{\parindent}{15pt} % Sangría estándar de párrafo (~0.5 cm)
\setlength{\parskip}{0pt}    % Sin salto extra entre párrafos estándar
\setstretch{1.25}            % Interlineado equivalente a 1.25 / 1.5 líneas (14.95 pt entre líneas base)

% =========================================================================
% FORMATO DE TÍTULOS DE SECCIÓN
% =========================================================================
% Numeración con punto final (1., 8.1., 8.1.1.)
\renewcommand{\thesection}{\arabic{section}.}
\renewcommand{\thesubsection}{\arabic{section}.\arabic{subsection}.}
\renewcommand{\thesubsubsection}{\arabic{section}.\arabic{subsection}.\arabic{subsubsection}.}

% Jerarquía de tamaños estándar LaTeX (14.35pt, 11.96pt, 9.96pt)
\titleformat{\section}{\Large\bfseries}{\thesection}{0.6em}{}
\titlespacing*{\section}{0pt}{18pt}{8pt}

\titleformat{\subsection}{\large\bfseries}{\thesubsection}{0.6em}{}
\titlespacing*{\subsection}{0pt}{14pt}{6pt}

\titleformat{\subsubsection}{\normalsize\bfseries}{\thesubsubsection}{0.6em}{}
\titlespacing*{\subsubsection}{0pt}{10pt}{4pt}

% =========================================================================
% METADATOS Y PARÁMETROS CONFIGURABLES DE LA PLANTILLA
% Reemplace los valores dentro de las llaves con los datos de su entregable.
% Parámetros estandarizados según Formato (1).pdf y Ejemplo formato entregables.pdf
% =========================================================================
\newcommand{\NombreOrganizacion}{[NOMBRE DE LA ORGANIZACIÓN / EMPRESA]}
\newcommand{\NombreProceso}{[NOMBRE DEL PROCESO]}
\newcommand{\TemaSubproceso}{[SUBPROCESO O ÁREA TEMÁTICA]}
\newcommand{\TituloDocumentoEncabezado}{[TÍTULO CORTO DEL DOCUMENTO]}

% Estructura oficial del código documental según Formato (1): TIPO_PROCESO-PROCESO-TIPO_DOC-SECUENCIAL
% Tipos de proceso: DI (Directivo), MI (Misional), AS (Asesor), AP (Apoyo)
% Procesos: TIC, DAF, SSO, ATH, PRE, TES, MAN, CPU, etc.
% Tipos doc: INF (Informe), PRO (Procedimiento), PLA (Planes), ESP (Especificaciones), etc.
\newcommand{\CodigoDoc}{XX-YYY-ZZZ-001}
\newcommand{\VersionDoc}{1.0}
\newcommand{\FechaElaboracion}{DD/MM/AAAA}
\newcommand{\FechaRevision}{DD/MM/AAAA}
\newcommand{\TituloCompleto}{[TÍTULO COMPLETO Y DESCRIPTIVO DEL INFORME / ENTREGABLE]}
\newcommand{\ArchivoLogo}{logo_devnova.png} % Logotipo institucional (3.0 cm x 3.0 cm)

% =========================================================================
% ESTILO DE PÁGINAS (ENCABEZADOS Y PIES)
% El documento de referencia mantiene las páginas limpias, sin encabezado
% superior, únicamente con el número de página centrado en el pie.
% =========================================================================
\pagestyle{plain}

% En caso de que en su carrera le soliciten encabezado formal en cada página:
% \pagestyle{fancy}
% \fancyhf{}
% \fancyhead[L]{\small\itshape\NombreOrganizacion}
% \fancyhead[R]{\small\itshape\CodigoDoc}
% \renewcommand{\headrulewidth}{0.4pt}
% \fancyfoot[C]{\thepage}

% =========================================================================
% INICIO DEL DOCUMENTO
% =========================================================================
\begin{document}

% -------------------------------------------------------------------------
% MEMBRETE INSTITUCIONAL OFICIAL (TABLA DE ENCABEZADO SEGÚN FORMATO 1)
% Tabla centrada de ancho fijo (~16.3 cm) con proporción oficial.
% -------------------------------------------------------------------------
\noindent
\begin{center}
\renewcommand{\arraystretch}{1.35}
\begin{tabular}{|>{\centering\arraybackslash}m{3.1cm}|>{\centering\arraybackslash}m{8.1cm}|>{\raggedright\arraybackslash}m{4.1cm}|}
\hline
\multirow{4}{3.1cm}{%
  \centering
  \IfFileExists{\ArchivoLogo}{%
    \includegraphics[width=3.0cm,height=3.0cm,keepaspectratio]{\ArchivoLogo}%
  }{%
    \fbox{\parbox[c][2.6cm][c]{2.8cm}{\centering\textbf{LOGO}}}%
  }%
}
& \textbf{\NombreOrganizacion} & \textbf{CÓDIGO:} \CodigoDoc \\ \cline{2-3}
& \textbf{PROCESO:} \NombreProceso & \textbf{VERSIÓN:} \VersionDoc \\ \cline{2-3}
& \TemaSubproceso & \textbf{ELABORACIÓN:} \FechaElaboracion \\ \cline{2-3}
& \TituloDocumentoEncabezado & \textbf{ÚLTIMA REVISIÓN:} \FechaRevision \\ \hline
\end{tabular}
\end{center}

\vspace{0.4cm}

% =========================================================================
% 1. TÍTULO
% =========================================================================
\section{Título}

\TituloCompleto

% =========================================================================
% 2. INTRODUCCIÓN
% =========================================================================
\section{Introducción}

[Describa aquí el contexto general del proyecto o área de trabajo, los antecedentes relevantes, la motivación u oportunidad abordada y una síntesis clara del propósito y contenido del informe. La introducción debe orientar al lector sobre la temática y los resultados que se presentan a continuación.]

[Puede continuar con un segundo párrafo describiendo la estructura del documento o las fuentes principales empleadas en el análisis.]

% =========================================================================
% 3. OBJETIVO
% =========================================================================
\section{Objetivo}

[Redacte el objetivo principal del documento comenzando con un verbo en infinitivo. Debe reflejar con precisión el propósito que se busca alcanzar con el informe o entregable, considerando la parte organizativa y técnica.]

% En caso de requerir desglose en objetivo general y específicos:
% \subsection{Objetivo General}
% [Redacte el objetivo general...]
%
% \subsection{Objetivos Específicos}
% \begin{itemize}
%     \item [Primer objetivo específico...]
%     \item [Segundo objetivo específico...]
%     \item [Tercer objetivo específico...]
% \end{itemize}

% =========================================================================
% 4. ALCANCE
% =========================================================================
\section{Alcance}

Este documento comprende:
\begin{itemize}
    \item La descripción y aplicación de los componentes, metodologías o normas abordadas.
    \item La integración de marcos de trabajo y herramientas técnicas utilizadas.
    \item La adopción de estándares de calidad, seguridad y gestión operativa.
    \item La incorporación de lineamientos institucionales y organizacionales del proyecto.
\end{itemize}

No se incluyen procesos administrativos externos ni evaluaciones fuera del marco temporal o técnico definido para el presente informe.

% =========================================================================
% 5. DEFINICIONES
% =========================================================================
\section{Definiciones}

\begin{itemize}
    \item \textbf{[Término o Concepto 1]:} [Definición clara y concisa en el contexto del proyecto o trabajo.]
    \item \textbf{[Término o Concepto 2]:} [Definición clara y concisa en el contexto del proyecto o trabajo.]
    \item \textbf{[Sigla o Acrónimo]:} [Significado de la sigla y breve explicación de su uso en el documento.]
\end{itemize}

% =========================================================================
% 6. RESPONSABILIDADES
% =========================================================================
\section{Responsabilidades}

[Detalle los miembros del equipo y sus responsabilidades directas relacionadas con las actividades o artefactos descritos en este documento.]

\begin{center}
\renewcommand{\arraystretch}{1.2}
\begin{tabular}{lp{10.5cm}}
\toprule
\textbf{Responsable} & \textbf{Descripción de la responsabilidad} \\
\midrule
{[Nombre y Apellidos 1]} & [Coordinar la planificación, aplicación y seguimiento general del entregable.] \\
{[Nombre y Apellidos 2]} & [Evaluar la infraestructura técnica y garantizar la aplicación de estándares técnicos.] \\
{[Nombre y Apellidos 3]} & [Supervisar la adopción de buenas prácticas de desarrollo y estándares de calidad.] \\
{[Nombre y Apellidos 4]} & [Documentar la implementación normativa y elaborar informes de cumplimiento.] \\
\bottomrule
\end{tabular}
\end{center}

% =========================================================================
% 7. NORMATIVA LEGAL Y TÉCNICA APLICABLE
% =========================================================================
\section{Normativa Legal y Técnica Aplicable}

[Liste las leyes, normativas institucionales y estándares técnicos internacionales que sustentan la elaboración del documento o rigen las actividades del proyecto.]

\begin{itemize}
    \item \textbf{[Normativa legal 1]:} [Descripción de la ley, decreto o código institucional aplicable.]
    \item \textbf{[Estándar Técnico 1 (ej. ISO/IEC 27001)]:} [Directrices para la seguridad de la información.]
    \item \textbf{[Estándar Técnico 2 (ej. ISO/IEC 25010)]:} [Criterios de evaluación de calidad del producto de software.]
    \item \textbf{[Marco de Referencia (ej. OWASP ASVS)]:} [Estándar de verificación de seguridad en aplicaciones.]
\end{itemize}

% =========================================================================
% 8. MÉTODO / POLÍTICAS / CONTENIDO
% =========================================================================
\section{Método / Políticas / Contenido}

\subsection{Metodología}

Se empleó un enfoque estructurado que combina análisis técnico y metodologías ágiles. La implementación se realizó aplicando principios de:

\begin{itemize}
    \item \textbf{Scrum:} con sprints definidos, revisión continua y roles asignados.
    \item \textbf{Gestión de Riesgos:} mediante matriz de riesgos y plan de contingencia.
    \item \textbf{Evaluación de Cumplimiento:} verificando conformidad con estándares técnicos.
\end{itemize}

\subsection{Políticas}

[Establezca las políticas, directrices o reglas obligatorias que rigen la ejecución de las actividades, la seguridad, la calidad o la toma de decisiones.]

\begin{itemize}
    \item \textbf{Política de Calidad:} [Directriz obligatoria sobre aseguramiento de calidad del entregable.]
    \item \textbf{Política de Seguridad:} [Directriz sobre protección de información y control de acceso.]
    \item \textbf{Política de Control de Versiones:} [Directriz sobre revisión y aprobación de artefactos.]
\end{itemize}

\subsection{Contenido}

[Desarrolle en esta sección el núcleo técnico del informe o entregable. Puede crear las subdivisiones temáticas que requiera mediante \texttt{\textbackslash subsubsection\{...\}}.]

\subsubsection{[Componente o Subtema Técnico 1]}

[Desarrollo detallado del componente, análisis, requerimientos, diseño o artefacto según el tipo de entregable.]

% -------------------------------------------------------------------------
% EJEMPLO FORMAL DE TABLA (Booktabs):
% -------------------------------------------------------------------------
% \begin{table}[htbp]
% \centering
% \caption{[Título descriptivo de la tabla o cuadro]}
% \label{tab:ejemplo}
% \renewcommand{\arraystretch}{1.2}
% \begin{tabular}{llcc}
% \toprule
% \textbf{Elemento} & \textbf{Descripción} & \textbf{Módulo} & \textbf{Estado} \\
% \midrule
% Módulo A & Procesamiento de autenticación & Backend & Conforme \\
% Módulo B & Interfaz de usuario interactiva & Frontend & En revisión \\
% \bottomrule
% \end{tabular}
% \end{table}

\subsubsection{[Componente o Subtema Técnico 2]}

[Continuación del desarrollo del contenido técnico.]

% =========================================================================
% 9. FIRMAS DE RESPONSABILIDAD
% =========================================================================
\section{Firmas de responsabilidad}

Indica la validez formal del documento y la identificación de las personas responsables de su elaboración, revisión y aprobación.

\begin{center}
\renewcommand{\arraystretch}{2.4}
\begin{tabular}{p{3.8cm}p{4.2cm}>{\centering\arraybackslash}p{4.0cm}>{\centering\arraybackslash}p{2.8cm}}
\toprule
\textbf{Rol} & \textbf{Nombre / Representante} & \textbf{Firma / Aprobación (QR)} & \textbf{Fecha} \\
\midrule
Gestor del proyecto & {[Nombre y Apellidos 1]} & & \FechaElaboracion \\
Desarrollador backend & {[Nombre y Apellidos 2]} & & \FechaElaboracion \\
Desarrollador frontend & {[Nombre y Apellidos 3]} & & \FechaElaboracion \\
Documentador & {[Nombre y Apellidos 4]} & & \FechaRevision \\
\bottomrule
\end{tabular}
\end{center}

% =========================================================================
% 10. CONTROL DE CAMBIOS
% =========================================================================
\section{Control de Cambios}

Al final del documento se registra el cuadro de control de cambios para mantener el historial de versiones y trazabilidad de revisiones.

\begin{table}[htbp]
\centering
\caption{Historial de control de cambios}
\label{tab:control_cambios}
\renewcommand{\arraystretch}{1.25}
\begin{tabular}{|c|p{8.0cm}|p{3.8cm}|c|}
\hline
\textbf{Versión} & \textbf{Descripción del Cambio} & \textbf{Responsable} & \textbf{Fecha} \\ \hline
1.0 & [Creación inicial y emisión del documento.] & {[Nombre del Responsable]} & \FechaElaboracion \\ \hline
% 1.1 & [Ajuste o modificación realizada] & {[Nombre del Responsable]} & [DD/MM/AAAA] \\ \hline
\end{tabular}
\end{table}

% =========================================================================
% 11. ANEXOS
% =========================================================================
\section{Anexos}

[Incluya aquí material de soporte, diagramas extensos, actas de reunión, especificaciones complementarias o evidencias que sirvan de sustento al informe.]

\appendix
\renewcommand{\thesection}{\Alph{section}.}

\section{[Nombre del Anexo o Documento de Soporte]}

[Descripción y contenido del Anexo A (diagramas de arquitectura, modelos de base de datos, actas, etc.).]

% -------------------------------------------------------------------------
% EJEMPLO DE FIGURA EN ANEXO:
% -------------------------------------------------------------------------
% \begin{figure}[htbp]
% \centering
% \includegraphics[width=0.85\textwidth]{nombre_imagen.png}
% \caption{Diagrama representativo del componente técnico.}
% \label{fig:diagrama_anexo}
% \end{figure}

\end{document}
